\documentclass[10pt,a4paper]{amsart}
\usepackage[margin=19mm]{geometry}
\usepackage{amsmath,amssymb,mathtools}
\usepackage[scr=rsfs,cal=euler]{mathalfa}
\usepackage{xfrac}
\usepackage[unicode,hidelinks,bookmarks=false]{hyperref}
\newcommand{\ie}{i.e.\ }
\newcommand{\cf}{cf.\ }
\newcommand{\resp}{respectively\ }
\newcommand{\Subsection}{\subsection}
\providecommand{\nobreakdash}{\nobreak\hbox{-}}
\setlength{\emergencystretch}{2em}
\multlinegap0pt
\usepackage{bm}
\usepackage{bbm}
\usepackage{dsfont}
\usepackage{xspace}
\usepackage{cancel}
\usepackage{mathtools}
\usepackage{amsthm}
\makeatletter
\@ifundefined{theorem}{\newtheorem{theorem}{Theorem}}{}
\@ifundefined{lemma}{\newtheorem{lemma}[theorem]{Lemma}}{}
\makeatother
\newcommand{\N}{\mathbb{N}}  
\newcommand{\R}{\mathbb{R}}
\title[The H\'enon equation in Orlicz-Sobolev spaces]{The H\'enon equation in Orlicz-Sobolev spaces}
\author{Pablo Ochoa, Ariel Salort}
\date{Source: arXiv:2509.17923v1 (CC BY 4.0)}
\begin{document}
\maketitle

\noindent\emph{Source and licence.} The H\'enon equation in Orlicz-Sobolev spaces. Version arXiv:2509.17923v1 (CC BY 4.0), distributed under the Creative Commons Attribution 4.0 International licence, as recorded in the arXiv licence metadata for this exact version and shown on the abstract page https://arxiv.org/abs/2509.17923v1. Full rights notice: licenses/arxiv-2509.17923-CC-BY-4.0.txt.\par\smallskip

\noindent\textit{Editorial scope.} Counted unit: the Lemma whose source label is \texttt{lemma.1} in the article (source0.tex lines 493--495, source label \texttt{lemma.1}), delivered together with the article's own complete original proof of it (source0.tex lines 496--542, one \texttt{proof} environment of 47 lines). No mathematics is written, repaired or re-derived: statement and proof are the article's own text line for line. Required context retained uncounted: none. Every definition, display and label of those lines is retained unchanged. Omitted from this package and not redistributed: 1--492, 543--1123. Nothing inside a retained excerpt is omitted or altered: every retained body is delivered exactly as the article's lines are written, so no omission receipt of any kind accompanies this package. Citations: the retained excerpts carry no citation command at all, so no bibliography entry of the article is transcribed and no reference is redistributed. Reference adaptation: none was needed and none was made; every reference inside the delivered unit resolves inside it, so no unresolved reference or citation is produced. Printed slips kept literal: no printed word, symbol, hypothesis or number of the counted statement or of its proof was altered. Layout and class: the article's own class is \texttt{amsart} with packages including amssymb, latexsym, esint, color, amsthm, amsfonts, xy, enumerate, mathrsfs, amsmath, verbatim, epsfig, graphicx, epstopdf; the wrapper is the lane's pinned \texttt{amsart} editorial wrapper at 10pt on A4, which restates the theorem-like environments the retained text needs and adds the article's own macro definitions that the retained text needs (\texttt{\textbackslash N}, \texttt{\textbackslash R}), with extra wrapper packages (bm, bbm, dsfont, xspace, cancel, mathtools). The delivered numbering is the unit's own. Assets: no retained span contains a figure environment, a graphics-inclusion command, a PDF-page inclusion, a table or a TikZ picture; no image file is copied into this package, the package declares no asset dependency and no image file is redistributed with it. Licence: version arXiv:2509.17923v1 of this article is distributed under the Creative Commons Attribution 4.0 International licence, as recorded in the arXiv licence metadata for this exact version and shown on the abstract page https://arxiv.org/abs/2509.17923v1. A full rights notice is delivered at licenses/arxiv-2509.17923-CC-BY-4.0.txt. Characters: every retained excerpt is pure ASCII, so the delivered engine, XeLaTeX with its default Latin Modern text font, reports no missing character.
\begin{lemma} \label{lemma.1}Let $\mu$ be a non-negative Radon measure absolutely continuous with respect to the Lebesgue measure in $\mathbb{R}^n$.
 Let $\Omega\subset \R^n$ be a bounded domain and let $\{u_k\}_{k \in \N}$ be a bounded sequence in $L_{\mu}^G(\Omega)$ such that $\{u_k\}_{k\in\N}$ converges to $u$ a.e. Then $u\in L_\mu^G(\Omega)$ and $u_k\to u$ strongly in $L_{\mu}^R(\Omega)$ for any $R\ll G$.
\end{lemma}

\begin{proof}
Let us observe that $u\in L_{\mu}^G(\Omega)$. Indeed, by Fatou's Lemma, we get that
$$
\int_\Omega G(|u|)\,d\mu \leq \liminf_{k\to\infty} \int_\Omega G(|u_k|)\,d\mu \leq C_0,
$$where we have used that $u_k$ converges to $u$ $\mu$-a.e. as well.

 
Let us prove that $u_k\to u$ strongly in $L_{\mu}^R(\Omega)$ when $R\ll G$. This condition means that 
$\frac{R(t)}{G(t)}\to 0$ as  $t\to\infty$. Then, for any $\varepsilon>0$, there is $T=T(\varepsilon)>0$ such that
\begin{equation}\label{R and G}
R(t)\leq \frac{\varepsilon}{8C} G(t) \quad \text{ for all } t \geq T,
\end{equation} 
with $C=C_0 c$, being $c$ the constant in the  $\Delta_2$-condition. For further reference, define for each $k>0$, the set:
$$\Omega_{T, k}:=\left\lbrace x\in \Omega: |u_k(x)-u(x)|\geq T \right\rbrace.$$

 Since $u_k\to u$ $\mu$-a.e. in $\Omega$, by Egorov's Theorem, given $\delta>0$ there exists a measurable set $E_\delta \subset \Omega$ such that $\mu(\Omega\setminus E_\delta)<\delta$ and $u_k\to u$ uniformly in $E_\delta$. Due to the uniform convergence, given $\delta>0$, there is $K>0$ such that $|u_k(x)-u(x)|\leq \delta$ for all $x\in E_\delta$ and for all $k>K$. Since $R$ is continuous and increasing, $R(|u_k-u|)\leq R(\delta)$ in $E_{\delta}$. We choose $\delta$ such that 
\begin{equation}\label{cond delta}
\delta<\min\left\lbrace T,\dfrac{ \varepsilon}{4R(T)}\right\rbrace, \quad R(\delta)\leq \frac{\varepsilon}{2\mu(\Omega)}
\end{equation}and take the corresponding $K=K(\delta)$. Then, for $k> K$, we get from \eqref{cond delta} that
$$
\int_{E_{\delta}} R(|u_k-u|)\,d\mu \leq R(\delta)\mu(\Omega)\leq \frac{\varepsilon}{2}.
$$

We next consider the set $\Omega\setminus E_{\delta}$.  Observe that  for $k> K$, we have 
$$
\Omega\setminus E_{\delta} = \left\lbrace x\in \Omega: |u_k(x)-u(x)|\geq T, \text{ for some }k\right\rbrace \cup \left\lbrace x\in \Omega: \delta \leq |u_k(x)-u(x)|< T, \text{ for some }k\right\rbrace,
$$
and moreover $\mu(\Omega\setminus E_{\delta})<\delta$. Observe that since $\delta<T$,
$$\Omega\setminus E_{\delta}= \Omega_{T, k} \cup \Omega \setminus (\Omega_{T, k}\cap E_{\delta}).$$  Then, for a fixed $k>K$ and recalling \eqref{R and G} and \eqref{cond delta}, we have
\begin{equation}
\begin{split}
\int_{\Omega\setminus E_\delta} R(|u_k-u|)\,d\mu& =   \int_{\Omega_{T, k}} R(|u_k-u|)\,d\mu + \int_{\Omega \setminus (\Omega_{T, k}\cap E_\delta)} R(|u_k-u|)\,d\mu\\ & \leq 
 \frac{\varepsilon}{8C}\int_{\Omega_{T, k}} G(|u_k-u|)\,d\mu + \int_{\Omega \setminus (\Omega_{T, k}\cap E_\delta)} R(T)\,d\mu\\ & \leq  \frac{c\varepsilon}{8C}\left( \int_{\Omega} G(|u_k|)\,d\mu  + \int_\Omega G(|u|)\,d\mu\right)+R(T)\delta
 \leq \frac{\varepsilon}{2}.
\end{split}
\end{equation}

\medskip



Therefore, for all $k>K$,
$$
\int_\Omega R(|u_k-u|)\,d\mu = \int_{\Omega\setminus E_\delta} R(|u_k-u|)\,d\mu + \int_{E_\delta} R(|u_k-u|)\,d\mu \leq \varepsilon
$$
which gives that $u_k\to u$ strongly in $L_{\mu}^R(\Omega)$.
\end{proof}

\end{document}
